\section{Related Work}
\label{sec:related}

\textbf{Representation Learning.} Representation learning aims to learn powerful semantic representations through diverse pretraining objectives. Contrastive methods such as SimCLR~\citep{simclr}, MoCo~\citep{moco,mocov2,mocov3}, and BYOL~\citep{byol} learn by maximizing agreement between augmented views of the same image. CLIP~\citep{clip} and SigLIP~\citep{siglip,siglip2} extend this paradigm to vision-language alignment, enabling zero-shot transfer across tasks. DINO variants~\citep{dino,dinov2,dinov3} train a student network to match a momentum-updated teacher. Masked autoencoding (MAE)~\citep{Vincent2010StackedDA,Pathak2016ContextEF} and its variants~\citep{mae,beit,ijepa} offer learning representations by reconstructing masked portions of the input.  

\textbf{Representation Alignment for Generation.} Existing methods can be categorized into those that align with external models and those that do not. Recent work has shown that aligning diffusion and flow model features with external pretrained encoders can significantly accelerate training and improve generation quality~\citep{repa,ldit,repae,repaet2i2025,wurstchen}, with extensions to domain-specific settings such as physics in video~\citep{videorepa} and geometry in 3D~\citep{Wu2025GeometryFM}. Another line of work trains generative models directly on semantic representations rather than reconstruction-driven latents~\citep{wu2025representation,rae}, though these methods remain tied to specific encoders which limits their reconstruction performance and hinders adaption across resolutions and modalities. For completeness, we show in Sec.~\ref{sec:experiments} and App.~\ref{sec:RAE} that our method improves the performance of RAE~\citep{rae}, demonstrating our method's robustness to autoencoder choice. 
As discussed in Sec.~\ref{sec:intro}, external alignment methods exhibit fundamental limitations such as unexpected scaling behavior (Sec.~\ref{sec:motivation}) and limited generalization across datasets and modalities (Sec.~\ref{sec:experiments}), motivating the need for a unified approach that eliminates dependence on external models entirely.

Existing methods that do not rely on external models can be broadly categorized into two groups. The first incorporates explicit self-supervised objectives at the cost of modifying the model's training dynamics, often necessitating an additional stage of pure diffusion fine-tuning to close the train-inference gap~\citep{maskdit,mdt,Chen2025MaskedAA,zhu2024sd}. The second preserves the diffusion framework~\citep{sra,dispersiveloss,haghighi2025layersyncselfaligningintermediatelayers} and employs diffusion features to perform alignment. However, the latter methods rely on the assumption that deeper layers naturally learn strong representations, while the former obtain results that are less favorable~\citep{dispersiveloss,maskdit,mdt}.
Overall, methods that use external representations have consistently outperformed those without. Our novel \emph{Self-Flow} approach closes this gap: by integrating self-supervised learning directly into flow matching, we surpass external alignment methods without requiring any external models.  